\section{Ring Switching}\label{sec:ring_switching}

Ring switching is a family of reductions that change the coefficient algebra used to represent
or check a claim. Coordinate packing appears in Binius \cite{DP24}, the generalized note
\cite{RSG}, and Flock \cite[Appendix B]{BRW26}. Hachi's subfield-to-cyclotomic trace reduction
\cite[\S3.1]{NOZ26} and the monic-quotient lift of \cite{HMZ25} have distinct relations and
security arguments. ArkLib separates their shared algebra from their protocol-specific checks,
challenge distributions, and commitment assumptions. No universal soundness loss follows from
the phrase ``ring switching.''

\subsection{Independent packing and evaluation algebras}

\begin{definition}[Finite-free packing data]
  \lean{RingSwitching.Packing.PackingData}
  \leanok
  \label{def:finite_free_packing_data}
  Let $B$ be a commutative ring. Choose independently two commutative $B$-algebras $P$ and $E$
  with finite bases $(\beta_i)_{i\in I}$ and $(\epsilon_u)_{u\in J}$. The indices may have
  different cardinalities. No embedding between $P$ and $E$, field structure, or domain
  hypothesis is required.
\end{definition}

\begin{definition}[Coordinate transpose]
  \lean{RingSwitching.Packing.PackingData.transpose}
  \leanok
  \uses{def:finite_free_packing_data}
  \label{def:packing_coordinate_transpose}
  The $B$-linear equivalence from $I\to E$ to $J\to P$ is
  \[
    T(\alpha)_u=\sum_i \beta_i\,[\alpha_i]_{\epsilon,u}.
  \]
  Its inverse unpacks each $P$-coordinate and recombines along the $E$-basis. Both inverse
  identities follow from the basis laws.
\end{definition}

\begin{definition}[Polynomial packing]
  \lean{RingSwitching.Packing.PackingData.polynomialEquiv}
  \leanok
  \uses{def:finite_free_packing_data}
  \label{def:packing_polynomial_equiv}
  A family of $B$-multilinears $(f_i)_{i\in I}$ in $m$ variables packs coefficientwise to
  $p=\sum_i\beta_i f_i$ over $P$. Unpacking each coefficient is the inverse operation.
  The equivalence also holds for $m=0$ and for commutative rings with zero divisors.
\end{definition}

\begin{theorem}[Family and slice read-back]
  \lean{RingSwitching.Packing.PackingData.openingClaimRel_iff_sliceRel}
  \leanok
  \uses{def:packing_coordinate_transpose, def:packing_polynomial_equiv}
  \label{thm:packing_family_slice}
  Fix $r\in E^m$, public values $(\alpha_i)$, and a packed polynomial $p$.
  The unpacked family satisfies $f_i(r)=\alpha_i$ for every $i$ if and only if
  \[
    T(\alpha)_u=
    \sum_{y\in\{0,1\}^m}[\widetilde{\mathrm{eq}}(r,y)]_{\epsilon,u}\,p(y)
    \quad\text{for every }u.
  \]
  This follows from the multilinear Boolean interpolation identity over commutative rings;
  no integral-domain interpolation assumption is needed.
\end{theorem}
\begin{proof}\leanok
Expand multilinears on the Boolean cube and apply the two inverse basis-coordinate maps.
\end{proof}

A compatible challenge algebra $C$ may extend $P$, with the same $B$-action along
$B\to P\to C$. For weights $w_u\in C$, define the $B$-linear map
$\Phi_w(z)=\sum_u[z]_{\epsilon,u}\,w_u$. The public multiplier is the multilinear extension
\[
  M_{r,w}=\operatorname{MLE}_C
    \bigl(y\mapsto\Phi_w(\widetilde{\mathrm{eq}}_E(r,y))\bigr).
\]
It is defined from a Boolean table. Since $\Phi_w$ need not preserve multiplication, its value
at an arbitrary $C$-point cannot be obtained by applying $\Phi_w$ to an unrelated $E$-evaluation.
The slice equations imply the sumcheck claim
\[
  \sum_u w_u\,\operatorname{map}_{P\to C}(T(\alpha)_u)
  =\sum_{y\in\{0,1\}^m}M_{r,w}(y)\,\operatorname{map}_{P\to C}(p(y)).
\]
Injectivity of $P\to C$ is a separate hypothesis for security read-back and separation; the
forward identity only needs compatible coefficient transport.

\begin{theorem}[Public multiplier evaluator]
  \lean{RingSwitching.Packing.PackingData.evaluateMultiplier_eq}
  \leanok
  \uses{def:finite_free_packing_data}
  \label{thm:packing_multiplier_evaluator}
  In the opening basis, form the multiplication matrices of $1-r_i$ and $r_i$ over $B$.
  Transport their entries into $C$ and interpolate layer $i$ as
  $(1-z_i)M_i(0)+z_iM_i(1)$. Applying these layers to the coordinates of $1$ and observing
  the final state with weights $w$ gives exactly $M_{r,w}(z)$. The observation is only linear;
  no embedding $E\to C$ or multiplicativity law is required.
\end{theorem}
\begin{proof}\leanok
Apply the read-once matrix interpolation identity to the actual opening-coordinate
multiplication matrices.
\end{proof}

The evaluator uses one matrix-vector action per retained variable, with width $|J|$, followed
by one weighted dot product. The instrumented recursion proves that action count; basis and
matrix preprocessing are excluded. This is a statement about the explicit arithmetic algorithm,
not an unproved runtime bound.

\subsection{Claim heads and batching}

The generalized note starts with the full public family $(\alpha_i)$ and derives its slices.
The implemented checked-slice presentation adds a prover message and checks that it equals
the same coordinate transpose before sampling a challenge. This is a distinct wire format.

Binius starts with one scalar evaluation. Its partial family must satisfy a checked equality
with Boolean equality weights. Flock's quirky claim head instead uses the concrete product of
Lagrange interpolation weights and Boolean equality weights. A zero-round conversion that merely
strengthens a scalar claim to the full-family relation is not a sound head.

\begin{theorem}[Scalar-source reconstruction]
  \lean{RingSwitching.Packing.ScalarHead.readback}
  \leanok
  \label{thm:packing_scalar_readback}
  A lossless component layout records the original scalar source, its retained evaluation point,
  and a proved weighted reconstruction. If the sent partial family passes that scalar check and
  satisfies the full-family relation on the same packed commitment, inverse component assembly
  recovers a source satisfying the original scalar claim. The actual one-message head has
  perfect completeness and zero-challenge knowledge soundness.
\end{theorem}
\begin{proof}\leanok
Use the layout reconstruction equation and inverse component assembly, preserving the packed
commitment.
\end{proof}

DP24's concrete layout packs the first Boolean block and retains the suffix; ordinary Flock
packs the suffix and retains the prefix. Their lossless Boolean-table splits hold over
commutative rings, including empty blocks.

\begin{theorem}[Flock quirky reconstruction]
  \lean{RingSwitching.Packing.ScalarHead.quirky_reconstruct}
  \leanok
  \label{thm:packing_quirky_readback}
  Let skipped indices enumerate distinct interpolation nodes in the opening field. Define the
  original quirky extension by Lagrange interpolation of the table's multilinear sections.
  Packing the pair $(\sigma,b)$ gives reconstruction weights
  $L_\sigma(\zeta)\widetilde{\mathrm{eq}}(\rho,b)$. The interpolant has degree below the number
  of nodes and agrees with the original section at each node.
\end{theorem}
\begin{proof}\leanok
Expand the Lagrange interpolation of the multilinear sections and apply the lossless component
split.
\end{proof}

Batching strategies carry a collision bound for two fixed distinct families. Over a finite
integral domain $C$, weights $1,c,\ldots,c^{e-1}$ give error at most
$(e-1)/|C|$ for $e\geq1$; equality folding over $\kappa$ coordinates gives
$\kappa/|C|$. Rank-one and empty-index cases have zero error. These root-count arguments do not
extend to arbitrary commutative rings by replacing the field cardinality with the ring cardinality.

A security event may quantify over witnesses after the challenge. Exact functionality of the
fixed commitment relation identifies them with a single polynomial chosen before that challenge,
which permits the fixed-family separation bound. The actual FRI-Binius initial compatibility
relation has proved functionality and honest coverage using codeword distance and the corrected
Boolean encoder. This premise does not prove the downstream dense opening. Flock's candidates
and out-of-domain selection, and Hachi's short-collision escape, need their own accounting.

\begin{definition}[Packed commitment relation]
  \lean{RingSwitching.Packing.PackedCommitment}
  \leanok
  \label{def:packing_commitment}
  An oracle family carries a fixed relation to packed $P$-polynomials and an honest constructor
  covering every polynomial. Its opening relation preserves the same commitment and asserts
  the polynomial's evaluation at a $C$-point. No uniqueness is assumed by this base interface.
\end{definition}

\begin{definition}[Exact packed commitment]
  \lean{RingSwitching.Packing.ExactPackedCommitment}
  \leanok
  \uses{def:packing_commitment}
  \label{def:packing_exact_commitment}
  The exact specialization additionally proves functionality for every fixed oracle collection.
  The generic phases, extractors, knowledge states and completeness use the base relation;
  randomized separation uses this functionality proof explicitly. Finite-list commitments can
  use the same phases, but require separate list/OOD probability bounds.
\end{definition}

\begin{theorem}[Checked full-family phase]
  \lean{RingSwitching.Packing.FullFamily.rbrKnowledgeSoundnessWorstCaseWith}
  \leanok
  \uses{thm:packing_family_slice, def:packing_exact_commitment}
  \label{thm:packing_full_family_knowledge}
  With compatible injective $P\to C$ transport and a proved batching strategy, the actual
  checked-slice verifier has worst-case-per-prefix knowledge error given by that strategy.
  Its extractor unpacks the committed polynomial. Failed read-back aborts, and the output
  retains the same oracle statements and commitment relation in the sumcheck claim.
\end{theorem}
\begin{proof}\leanok
Functionality fixes the committed polynomial before sampling; injective transport and the
batching separation bound control the bad event.
\end{proof}

\begin{theorem}[Actual scalar and family composition]
  \lean{RingSwitching.Packing.ScalarFamily.rbrKnowledgeSoundnessWorstCaseWith}
  \leanok
  \uses{thm:packing_scalar_readback, thm:packing_full_family_knowledge}
  \label{thm:packing_scalar_family_knowledge}
  Appending the actual checked scalar head and checked family phase preserves both guards and
  the exact composed extractor and knowledge state. Under the family phase's explicit security
  hypotheses, the only challenge has its batching error. Stateful perfect completeness also
  holds for the base commitment relation, including a genuinely ambiguous finite candidate set.
  The composition sends partial values and redundant checked slices in two consecutive messages,
  whereas the DP24 and Flock source protocols send a single partial-family message.
\end{theorem}
\begin{proof}\leanok
Apply guarded-first knowledge composition to the proved zero-error scalar head and checked
family phase.
\end{proof}

The concrete Binius adapter supplies this commitment contract from the production first-codeword
relation. GF(16) clients instantiate the actual phase and complete generic sumcheck pipeline,
including honest source coverage and rejection while holding the commitment fixed. FRI-Binius's
subsequent protocol interleaves FRI and sumcheck; it is not the legacy sequential dense-opening
wrapper below.

\subsection{The generic product-sumcheck tail}

For a public context selecting a multilinear multiplier $M\in C[X_1,\ldots,X_m]$, the tail
tracks the residual Boolean sum of $M\cdot\operatorname{map}_{P\to C}(p)$ after a prefix of
scalar challenges. Each factor is multilinear, so the product has individual degree at most two.
The original $P$-polynomial and commitment oracle family remain in every relation.

\begin{theorem}[Generic sumcheck round]
  \lean{RingSwitching.Packing.Tail.Round.rbrKnowledgeSoundnessWorstCaseWith}
  \leanok
  \uses{def:packing_commitment}
  \label{thm:packing_tail_round}
  The actual verifier reads a degree-two polynomial $g$, checks $g(0)+g(1)$ against its input
  target, and extends the prefix by a uniform scalar challenge $c\in C$ with target $g(c)$.
  Before the challenge, knowledge requires the genuine residual polynomial. Afterward it
  requires the local guard and the actual sampled residual relation. Over a finite integral
  domain, explicit commitment functionality fixes the witness before sampling and gives
  worst-case-per-prefix error $2/|C|$. Honest execution is perfectly complete over finite
  commutative challenge rings.
\end{theorem}
\begin{proof}\leanok
Fix the committed polynomial before sampling and bound roots of the degree-two difference from
the genuine residual polynomial.
\end{proof}

\begin{theorem}[Generic terminal opening]
  \lean{RingSwitching.Packing.Tail.Terminal.rbrKnowledgeSoundnessWorstCaseWith}
  \leanok
  \uses{def:packing_commitment}
  \label{thm:packing_tail_terminal}
  After all variables are fixed, the prover sends $v$. The verifier checks
  $s=M(r')v$ and forwards $(r',v)$ on the same commitment. Correct output read-back needs
  no cancellation and no functionality; the deterministic step adds no challenge error.
  With $m=0$, the empty sequence is identity and reaches this same terminal message.
\end{theorem}
\begin{proof}\leanok
Substitute the accepted packed evaluation into the local guard and the terminal residual
identity; no multiplier cancellation is used.
\end{proof}

\begin{theorem}[Actual generic packing pipelines]
  \lean{RingSwitching.Packing.FullFamilyOpening.rbrKnowledgeSoundnessWorstCaseWith,
        RingSwitching.Packing.ScalarOpening.rbrKnowledgeSoundnessWorstCaseWith}
  \leanok
  \uses{thm:packing_full_family_knowledge, thm:packing_scalar_family_knowledge,
        thm:packing_tail_round, thm:packing_tail_terminal}
  \label{thm:packing_generic_pipeline}
  Actual sequential composition of the corresponding head, the $m$ scalar rounds and terminal
  check reduces the original full-family or scalar relation to the same commitment's
  $C$-valued opening relation. The exact extractor and knowledge state use the proved guarded
  composition constructors. Under explicit functionality, finite-domain challenges and
  compatible injective $P\to C$ transport, the batching challenge retains its strategy error
  and each tail challenge has error $2/|C|$. Averaging preserves these exact objects.
  Stateful perfect completeness uses the base commitment relation without functionality.
\end{theorem}
\begin{proof}\leanok
Compose the actual guarded head, finite sequence of product-sumcheck rounds, and terminal check
using their exact extractors and knowledge states.
\end{proof}

\begin{theorem}[Exact challenge-error accounting]
  \lean{RingSwitching.Packing.FullFamilyOpening.rbrError_sum,
        RingSwitching.Packing.ScalarOpening.rbrError_sum}
  \leanok
  \uses{thm:packing_generic_pipeline}
  \label{thm:packing_error_sum}
  The sum of the declared errors over the actual challenge indices is
  $\varepsilon_{\mathrm{batch}}+m(2/|C|)$ in both pipelines. The deterministic scalar
  head, empty adapter and terminal message add no challenge. This arithmetic identity
  does not assert an ordinary-soundness bound.
\end{theorem}
\begin{proof}\leanok
Use the canonical sequence and append challenge-index equivalences and the public component error
lookup identities.
\end{proof}

\begin{definition}[Same-commitment downstream opening]
  \lean{RingSwitching.Packing.PackedOpening}
  \leanok
  \uses{def:packing_commitment}
  \label{def:packing_downstream_opening}
  A downstream opening argument is an actual oracle reduction whose input relation is precisely
  the packed commitment's $C$-valued evaluation relation. Its output statement, oracle family,
  witness and relation are explicit. The front and opening share the same ambient oracle.
\end{definition}

\begin{theorem}[Actual downstream assembly]
  \lean{RingSwitching.Packing.PackedOpening.rbrKnowledgeSoundnessWorstCaseWith,
        RingSwitching.Packing.PackedOpening.perfectCompleteness}
  \leanok
  \uses{thm:packing_generic_pipeline, def:packing_downstream_opening}
  \label{thm:packing_downstream_assembly}
  Actual append composes supplied worst-case knowledge contracts at the same evaluation seam,
  using their exact extractors and guarded knowledge states. The downstream verifier may have
  arbitrary effects. Perfect completeness separately requires guarded verification, completeness
  from every seam state, and the stated pure-output or first-message condition.
\end{theorem}
\begin{proof}\leanok
Apply guarded-first knowledge append and, separately, guarded completeness append with the
stated seam-state and execution hypotheses.
\end{proof}

The concrete packing fronts use the empty ambient oracle. Checked polynomial-oracle clients
close both public pipelines and prove probability-one acceptance of a nonconstant source.
A separate generic front with nonempty ambient oracle verifies that an accepted front executes
the downstream state-changing query and that a rejected front prevents it.
A downstream PCS still supplies its own proved security contract; the admitted unrestricted
composition and ordinary-security bridges are not used here.

\subsection{The legacy tensor specialization}

\begin{definition}[Faithful tensor profile]
  \lean{RingSwitching.RingSwitchingProfile}
  \leanok
  \label{def:ring_switching_profile}
  The legacy DP24 pipeline fixes one finite-free $B$-algebra $L$, a carrier $A$, two embeddings
  $\varphi_0,\varphi_1:L\to A$, and row and column coordinates. Each coordinate map is a
  two-sided inverse to its corresponding recomposition:
  \[
    R(c)=\sum_u\varphi_0(c_u)\varphi_1(\beta_u),\qquad
    Q(c)=\sum_u\varphi_1(c_u)\varphi_0(\beta_u).
  \]
  The embeddings agree on $B$. One-sided reconstruction alone is insufficient: an affine or
  collapsed coordinate map can satisfy it while rejecting an honest evaluation claim.
\end{definition}

\begin{definition}[Binius tensor instance]
  \lean{RingSwitching.binaryTowerProfile}
  \leanok
  \uses{def:ring_switching_profile}
  \label{def:binary_tower_profile}
  Set $A=L\otimes_B L$, $\varphi_0(x)=x\otimes1$, and $\varphi_1(y)=1\otimes y$.
  The left and right base-changed bases provide the respective coordinate maps and prove all
  profile laws. This algebraic constructor accepts commutative rings.
\end{definition}

The pure-tensor laws determine the orientation:
\[
  \operatorname{Rows}(\varphi_0(x)\varphi_1(y))_u=x[y]_{\beta,u},\qquad
  \operatorname{Columns}(\varphi_0(x)\varphi_1(y))_u=y[x]_{\beta,u}.
\]
Thus rows recover the original partial evaluations, while columns retain packed values for
batching and the final multiplier. Permanent production tests include an honest nonzero source
that the former reversed coordinate check rejected.

\begin{definition}[Boolean-table packing]
  \lean{RingSwitching.packMLE}
  \leanok
  \uses{def:ring_switching_profile}
  \label{def:pack_mle}
  For $\ell=\kappa+\ell'$, the packed polynomial has Boolean values
  \[
    t'(w)=\sum_{v\in\{0,1\}^{\kappa}}\beta_v\,t(v,w).
  \]
  This groups Boolean table entries. Hachi's monomial-coefficient convention is a different
  layout and requires its own correspondence.
\end{definition}

\begin{definition}[Legacy full reduction]
  \lean{RingSwitching.FullRingSwitching.fullOracleReduction}
  \leanok
  \uses{def:ring_switching_profile, def:pack_mle}
  \label{def:full_ring_switching}
  The legacy construction composes its checked scalar/batching phase, degree-two sumcheck,
  and a supplied dense multilinear opening protocol.
\end{definition}

Every failed local check now causes absorbing monadic failure. At the terminal check
$\mathrm{target}=M(r')s'$, the verifier forwards the opening value $s'$, including when
$M(r')=0$. The final message has no challenge index and contributes no probabilistic
root-count error. A sumcheck collision at one sampled point does not imply that the preceding
sum target was true or that two whole polynomials were equal.

The final leaf has proved ring-valid completeness and zero-error knowledge soundness, including
its exact execution and structural witness reconstruction. The remaining legacy batching/loop
completeness and knowledge leaves are admitted. They are proof obligations,
not certificates of end-to-end Binius security. The proved guarded-first knowledge-composition
specialization accepts worst-case-per-prefix component bounds and preserves their exact extractors;
an averaged-only PCS bound cannot be supplied as a worst-case premise. State-aware guarded
completeness composition likewise requires the appropriate suffix-state contract.

\subsection{Trace relocation and quotient lifting}

Hachi's deterministic head uses the monomial packing $\psi$ and subfield-valued points. Its
trace pairing is scaled:
\[
  \operatorname{Tr}_H(\psi(a)\sigma_{-1}(\psi(b)))=(d/k)\langle a,b\rangle.
\]
\begin{theorem}[Exact scalar trace check]
  \lean{ArkLib.Lattices.Hachi.TraceHead.trace_eval_eq_iff}
  \leanok
  \label{thm:hachi_scalar_trace_check}
  Pack the last variables' monomial coefficients using the actual $\psi$ equivalence.
  At retained points embedded from the fixed subring, the scaled trace check on the packed
  polynomial's ring evaluation holds if and only if its decoded scalar polynomial has the
  original claimed evaluation. The factor $d/k$ is a unit in the cyclotomic ring by odd
  characteristic and the power-of-two rank. No fixed-subring field instance is used.
\end{theorem}
\begin{proof}\leanok
Apply the scaled trace pairing to coefficient packing and cancel its proved unit factor.
\end{proof}

\begin{theorem}[Honest scalar commitment coverage]
  \lean{ArkLib.Lattices.Hachi.TraceHead.committed_source_valid}
  \leanok
  \uses{thm:hachi_scalar_trace_check}
  \label{thm:hachi_scalar_commitment_coverage}
  For every scalar coefficient polynomial and query, apply the existing Hachi committer to
  its coefficient packing. Balanced gadget reconstruction recovers that same ring polynomial
  from the honest decommitment. Under the existing digit/norm inequalities, it supplies a valid
  source weak opening, including the message bound required by the nonrecursive honest chain.
\end{theorem}
\begin{proof}\leanok
Use balanced-gadget reconstruction of the committed packing and the existing honest-opening norm
bounds.
\end{proof}

The actual one-message head has perfect completeness and zero-challenge CWSS, preserving the
same commitment and norm-conditioned weak opening at the existing ring-evaluation relation.
Its semantic definitions use the existing noncomputable fixed-subring and trace infrastructure.
Packaging an executable scalar scheme, the external field identification, and downstream CWSS
collision accounting are separate obligations.

The generic quotient lift in \texttt{RingSwitching/Lift/} instead moves an identity modulo a
monic polynomial to an exact lifted polynomial identity, followed by evaluation and interpolation
in a field. Its cyclotomic instance is Hachi \S4.3, already part of the nonrecursive chain.
Exceptional-set interpolation over wider Galois-ring targets remains separate work. Hachi's
\S3.2/\S4.5 recursion compression also needs a protocol repair with explicit rank, layout,
shortness, and commitment correspondence; the packing algebra does not supply it.
